% !TeX root = main_long.tex
\section{Introduction}

Healthcare research and personalized medicine depend on integrating clinical observations, molecular data, and biomedical knowledge across institutions and databases.
Differences in data structures, identifiers, and terminology complicate their combined use~\cite{gaudetblavignac2021sphn,putman2024monarch}.
Semantic Web approaches address this heterogeneity through shared identifiers, explicit relationships, and reusable vocabularies~\cite{smith2007obo,belleau2008bio2rdf}, but meaningful integration also depends on preserving the context in which observations were recorded.

The Variant Call Format (VCF) is widely used to exchange genomic variant data~\cite{danecek2011vcf}.
Each record describes a genomic location, reference and alternative alleles, and optional annotations.
Sample columns carry observations such as genotype and sequencing depth.
Header declarations define fields and how their values should be interpreted~\cite{vcf45}.
The GA4GH-maintained specifications define this conserved structure across VCF 4.1--4.5, while differing in declarations and constraints~\cite{ga4gh_hts_specs,htsspecs}.
VCF content fields can contain dependencies to other fields, essentiually encoding meaning through the structure of the file.
For example, genotype \term{0/2} identifies the reference allele and the second listed alternative allele, which is indexed based on the order in which these alternate alleles appear in the file, while separately an annotation's cardinality may depend on its header declaration.
When observations are integrated with external information, those relationships must remain available.

RDF and Linked Data offer a basis for exposing this context as explicit, addressable relationships alongside links to biomedical resources~\cite{decker2000semantic,bizer2009linked,lubin2017variantfile}.
This supports the interoperability and reuse sought by the FAIR principles~\cite{wilkinson2016fair}, although wider FAIRness also depends on identifiers, metadata, access, and provenance~\cite{bahim2020fair}.
Our aim is to provide a reusable representation of VCF observations that retains their source context to better enable these connections.

\subsection{Related Work}
Related work includes reusable semantic models and VCF--RDF conversion or query tools, with some contributions combining both roles (Table~\ref{tab:vcf-rdf-related-work}).
The Genomic Feature and Variation Ontology (GFVO), Genome Variation Ontology (GVO), and HERO-Genomics model genomic or clinical concepts.
Tools such as VCF2RDF, TogoVar, BioInterchange, AgroLD, and bio-vcf~\cite{garrison2022vcflib} expose VCF-derived data for integration.
The GA4GH Variation Representation Specification (VRS) complements these approaches with computable variation descriptions and federated identification, using JSON Schema rather than RDF or OWL; it does not preserve source-file structure~\cite{wagner2021vrs}.

These approaches distinguish source observations from biological interpretations.
VCF Core provides a shared, version-aware representation of source observations with links to existing genomic models, allowing converters and queries to reuse VCF interpretation rules.

\begin{table}[htbp]
\caption{VCF-related approaches and their roles: reusable models, RDF conversion/query tools, or combinations of both.}
\label{tab:vcf-rdf-related-work}
\centering
\small
\begin{tabularx}{\linewidth}{@{}>{\raggedright\arraybackslash}p{0.27\linewidth}>{\raggedright\arraybackslash}X@{}}
\toprule
Approach & Role and focus \\
\midrule
BioHackathon \textit{vcf2rdf}~\cite{katayama2019biohackathon}
& \textbf{Converter.} Early VCF--RDF conversion, using FALDO for positions and tool-specific terms for other VCF concepts.
\\
GFVO~\cite{baran2015gfvo}
& \textbf{OWL ontology.} Harmonizes genomic features and variation across VCF, GVF, GFF3, and GTF.
\\
BioInterchange~\cite{baran2015gfvo}
& \textbf{Converter.} Emits GFVO from GFF3, GVF, and VCF; its VCF reader extends the GFF3 reader, so records become generic features, and it reads VCF 4.1--4.2 only.
\\
VCF2RDF~\cite{penha2017vcf2rdf}
& \textbf{Mapping/converter.} Isomorphic VCF--RDF mapping exposing file components independently with source context; not a genomic ontology.
\\
TogoVar \textit{vcf2rdf}~\cite{mitsuhashi2022togovar}
& \textbf{Converter.} Simple RDF representation of population allele-frequency VCF data for integration with TogoVar's biomedical knowledge graphs.
\\
GVO~\cite{kawashima2023gvo}
& \textbf{OWL ontology.} Describes biological variation, particularly complex structural variants, rather than VCF structure and interpretation rules.
\\
VariantKG~\cite{prasanna2025variantkg}
& \textbf{Converter + application ontology.} Builds VCF-derived graphs for variant analysis and graph machine learning, not specification modelling.
\\
SPARQLing genomics~\cite{janssen2020sparqlinggenomics}
& \textbf{Converter + query platform.} Generates RDF from VCF and other genomic sources and serves it for querying; its terms are tool-specific, not a VCF model.
\\
Semantic Beacons~\cite{bodrug2025semantic}
& \textbf{Integration mapping.} Maps GA4GH Beacon responses to RDF with FALDO, SO and GENO for federated genomic querying; not a VCF--RDF representation.
\\
HERO-Genomics~\cite{cazzaro2025hero}
& \textbf{Ontology.} Integrates VCF- and Beacon-derived information across clinical and research settings, with broader scope than VCF specification coverage.
\\
AgroLD VCF2RDF~\cite{larmande2025agrold}
& \textbf{Converter.} Incorporates VCF variation into AgroLD's plant-science knowledge graph; not a dedicated VCF ontology.
\\
VCF-RDFizer~\cite{crum2025semantifying}
& \textbf{Converter + precursor vocabulary.} Our earlier VCF--RDF workflow; VCF Core separates model from conversion and narrows scope to the specification.
\\
\bottomrule
\end{tabularx}
\end{table}

\paragraph{The VCF Core Vocabulary.}
Here, we present and assess the \href{https://ecrum19.github.io/vcf-core-vocabulary/}{VCF Core Vocabulary}\footnote{\url{https://ecrum19.github.io/vcf-core-vocabulary/}} (\textit{vcfcv})~\cite{vcfc210}.
It combines source-context preservation, \textit{expanded} and \textit{condensed} sample column representations, 123 curated alignments to existing semantic models, and version-specific artifacts for VCF 4.1--4.5.
Four questions structure its assessment:
\textbf{RQ1:} does VCF Core preserve the intended meaning of the VCF constructs it claims to support?
\textbf{RQ2:} do the version-specific artifacts, the two \emph{sample profiles}, and a semantic round-trip behave as documented in graphs from a separately implemented VCF-RDF converter?
\textbf{RQ3:} does the resulting representation support a reproducible integration task while retaining source context?
\textbf{RQ4:} how does VCF Core relate to other existing VCF semantic models?

We assess representation, implementation, and answer correctness separately, combining specification-derived tests, comparison across producers, an integration example, and comparison with related models.
The resulting evidence identifies both reusable capabilities and gaps requiring further testing.
Throughout, \term{vcfc:} denotes \url{https://w3id.org/vcf-core/vocab#}.
